\appendix
\phantomsection
\addcontentsline{toc}{chapter}{Appendix}

\chapter{Key Simulation Code}
\label{app:key-code}

This appendix retains only condensed algorithm extracts for the rules used
in Chapters~6 and~7. Initialization, plotting, file export, diagnostics, and
Monte Carlo orchestration are omitted. The listings provide a formula-to-code
map: centralized matching every
business day (\texttt{centralized\_cycle\_length=1}); RFQ with four rounds
and \(K=4\); installment origination when rollover is ON and next-day single
payment when it is OFF; and early stopping only when one or fewer
non-central banks remain active. The complete executable source, realized
seeds, raw artifacts, dependency specification, and figure/table scripts are
supplied in the electronic reproduction materials described in Section~6.1.
The study's repository address is
\begin{center}\small
\url{https://github.com/zhanghaoyu1232-debug/central-bank-simulation}
\end{center}
The supplied core source and formal artifacts share fingerprint
\texttt{11dda58061ded609}; the reproduction manifest records full-file
checksums. The listings below are condensed extracts; the complete
source remains the executable reference.

\section{Contracts, Origination-Based Rollover, Gross Due-Payment Aggregation, and Role Assignment}

Listing~\ref{lst:model-core} reflects the current contract treatment. The
rollover switch selects the repayment schedule when the trade is originated.
Rollover ON registers a 5--20-period amortizing installment contract
immediately. Rollover OFF registers a one-business-day single-payment
contract. The current code does not first create a 10-day bullet loan and does
not convert principal at a later maturity date. The listing also shows the
gross debtor-to-creditor orientation used to build due-payment liabilities.

\lstinputlisting[
  caption={Current contract structures, origination-based rollover, gross due-payment aggregation, and role assignment},
  label={lst:model-core}
]{backmatter/code/model_core.py}

\section{Centralized Matching}

The centralized mechanism updates each bank's latest order and executes a
market-wide allocation on every business day. The formal experiments set the
matching cycle length to 1, so there is no 15-business-day batch delay.
Executable quantities are reconfirmed from current bank states. The allocation
is rate-prioritized and subject to capacity, degree, and single-loan
constraints. Borrower funding is all-or-nothing: provisional allocations are
committed only when the complete demand can be assembled.

\lstinputlisting[
  caption={Daily centralized matching with a borrower-level full-funding rule},
  label={lst:centralized-core}
]{backmatter/code/centralized_matching_core.py}

\section{Decentralized RFQ Matching}

The decentralized RFQ mechanism is executed every business day. Each borrower
constructs one local lender pool. Outstanding rollover lenders are retained
and do not consume the new-discovery budget of eight lenders. The borrower may
then conduct up to four quotation rounds and contact at most \(K=4\) lenders
per round. Accepted bilateral trades are retained even when the borrower's
total demand is not fully satisfied. The listing records every lender source,
the discovery and random-exploration limits, exclusion after rejection or
capacity exhaustion, the changing candidates across rounds, and the retained
partial-funding rule. The complete GNN feature and scoring functions are
included in the electronic source snapshot.

\lstinputlisting[
  caption={Four-round decentralized RFQ matching with local counterparty search and retained partial funding},
  label={lst:rfq-core}
]{backmatter/code/decentralized_rfq_core.py}

\section{Eisenberg--Noe Clearing, Creditor Losses, and Systemic Risk}

Listing~\ref{lst:clearing-risk-core} uses a non-negative debtor-to-creditor
liability matrix, so reciprocal obligations remain separate rather than being
bilaterally netted. Actual creditor receipts are determined by the
Eisenberg--Noe fixed point, and bilateral unpaid claims are then aggregated by
creditor. Accounting equity is evaluated before regulatory
core capital is truncated at zero, allowing negative-equity insolvency to be
identified explicitly. No contemporaneous liquidity facility is inserted into
the clearing matrix.

\lstinputlisting[
  caption={Gross Eisenberg--Noe clearing, creditor-payment allocation, accounting-equity failure, and systemic-risk measurement},
  label={lst:clearing-risk-core}
]{backmatter/code/clearing_and_risk_core.py}

\section{Period Update Sequence}

Listing~\ref{lst:period-update-core} summarizes the current business-day
sequence. Queued liquidity support is disbursed first. Currently due interbank
payments are then settled, after which eligible liquidity support is queued
for the next business day. Matching forms new installment or single-payment
contracts according to the rollover switch. After project returns,
solvency/capital support may be applied. Unpaid ON installments become
arrears, unpaid OFF single payments cause hard default, and creditor
write-offs are propagated recursively until no new insolvency appears. The
post-EN negative-equity cascade is entered even when the payment-default
list is empty. Each cascade iteration rebuilds positions, calculates
untruncated accounting equity, and resolves newly insolvent banks before
the next iteration. The post-project cascade is executed before risk
measurement. The adjacent resolution extract shows the creditor recovery
pool and write-off calculation. Network stability is recorded as a
diagnostic and does not terminate a formal run. A run ends at day
\(T=1000\) unless one or fewer non-central banks remain active.

\lstinputlisting[
  caption={Condensed current period sequence with daily matching, origination-based rollover, and horizon-based stopping},
  label={lst:period-update-core}
]{backmatter/code/period_update_core.py}

The resolution batch freezes the failure cohort before distributing
recoveries. Claims held by a failed lender against surviving borrowers
are transferred to the central bank at their recovery value, subject to
central-bank cash. The resulting estate cash is distributed pro rata to
eligible surviving creditors, up to the LGD recovery cap. Original claims
are removed, interbank positions are rebuilt, and equity is recomputed;
the diagnostic write-off is not deducted a second time from equity.

\lstinputlisting[
 caption={Creditor recovery and write-offs within the absorbing-default batch},
 label={lst:resolution-core}
]{backmatter/code/resolution_core.py}

\chapter{Initial Bank Data for Representative Runs}
\label{app:initial-bank-data}
% Legacy label retained so older cross-references do not break.
\label{app:policy-event-logs}

This appendix reports the shared initial-state export accompanying the
paired experiments in Chapters~6 and~7. The centralized and decentralized
CSV snapshots agree across all 30 banks and all 22 fields and are printed
once. They illustrate the initialization; individual Monte Carlo states
are regenerated from the realized replication seeds, rather than loading
this same CSV as a fixed input to every replication. The tables contain
two balance-sheet groups and two regulatory/risk groups.

The exported solvency ratio is accounting equity divided by total
liabilities, as in Equation~\eqref{eq:solvency}. The initial LCR field is
a construction-time cache calculated with outflow 0.4 for non-central
banks and 0.2 for the central bank; the separately exported outflow field
is the bank's randomized initial rate. At regulatory refresh, LCR is
recomputed using the current rate and Equation~\eqref{eq:lcr}. Thus, for
Bank1, the initial solvency ratio is 0.5113 and the cached LCR is 1.2623.
These values retain the recorded initialization rather than imposing a
later-period recomputation on the exported snapshot.

\begin{table}[htbp]
\centering
\caption{Overview of the shared initial-state export}
\label{tab:initial-bank-data-overview}
\small
\begin{tabular}{lrrrrl}
\toprule
Dataset & Banks & Central & Commercial & Shadow & Role in the thesis \\
\midrule
Shared initial export & 30 & 1 & 20 & 9 & Representative snapshot \\
\bottomrule
\end{tabular}
\end{table}

All monetary values are shown in simulation units and rounded to two
decimal places. Solvency, capital, liquidity, leverage, risk-appetite,
volatility, and outflow variables are shown to four decimal places; loan
and investment rates are shown to six decimal places. ``IB'' denotes
interbank.

\begingroup
\small
\setlength{\tabcolsep}{5pt}
\renewcommand{\arraystretch}{1.12}

\section{Shared Initial Snapshot for Centralized Matching and Decentralized RFQ}

\begin{longtable}{@{}llrrrr@{}}
\caption{Shared initial data: core capital, liquid assets, current liabilities, and deposits}\label{tab:paired-current-initial-balance-sheet}\\
\toprule
Bank & Type & Core capital & Liquid assets & Current liab. & Deposits \\
\midrule
\endfirsthead
\multicolumn{6}{c}{\tablename\ \thetable\ -- continued}\\
\toprule
Bank & Type & Core capital & Liquid assets & Current liab. & Deposits \\
\midrule
\endhead
\midrule
\multicolumn{6}{r}{Continued on the next page}\\
\endfoot
\bottomrule
\endlastfoot
CentralBank & Central & 5445.20 & 5388.65 & 3886.54 & 3095.81 \\
Bank1 & Commercial & 1640.55 & 1620.20 & 3208.84 & 2454.11 \\
Bank2 & Commercial & 608.48 & 1028.73 & 4202.48 & 3084.42 \\
Bank3 & Commercial & 2850.49 & 2649.87 & 4342.39 & 3037.61 \\
Bank4 & Commercial & 755.66 & 878.97 & 1233.04 & 853.94 \\
Bank5 & Commercial & 1787.89 & 2052.02 & 2641.36 & 1883.37 \\
Bank6 & Commercial & 1744.95 & 2072.82 & 3278.73 & 2337.29 \\
Bank7 & Commercial & 1851.98 & 2181.60 & 3296.21 & 2435.22 \\
Bank8 & Commercial & 720.78 & 1119.80 & 3990.22 & 2932.71 \\
Bank9 & Commercial & 2420.66 & 2358.71 & 3981.77 & 2839.30 \\
Bank10 & Commercial & 1998.03 & 2241.06 & 2430.33 & 1762.10 \\
Bank11 & Commercial & 1906.01 & 2117.04 & 2110.30 & 1575.95 \\
Bank12 & Commercial & 2064.19 & 2471.26 & 4070.66 & 2857.75 \\
Bank13 & Commercial & 929.86 & 1243.04 & 3131.73 & 2227.72 \\
Bank14 & Commercial & 1491.41 & 1911.18 & 4197.68 & 3082.29 \\
Bank15 & Commercial & 1057.50 & 1194.59 & 3404.16 & 2240.01 \\
Bank16 & Commercial & 1035.67 & 1397.32 & 3616.51 & 2856.07 \\
Bank17 & Commercial & 1511.69 & 1740.81 & 2291.17 & 1731.08 \\
Bank18 & Commercial & 602.45 & 730.28 & 1278.27 & 952.91 \\
Bank19 & Commercial & 920.91 & 1133.85 & 2129.48 & 1456.06 \\
Bank20 & Commercial & 666.80 & 881.95 & 2151.53 & 1508.90 \\
Bank21 & Shadow & 1710.74 & 2015.27 & 3045.25 & 469.87 \\
Bank22 & Shadow & 1116.76 & 836.31 & 2030.12 & 321.77 \\
Bank23 & Shadow & 682.44 & 825.58 & 1431.33 & 219.40 \\
Bank24 & Shadow & 746.91 & 982.87 & 2359.66 & 363.49 \\
Bank25 & Shadow & 1946.41 & 2036.40 & 899.89 & 128.57 \\
Bank26 & Shadow & 1080.97 & 1395.42 & 3144.42 & 477.40 \\
Bank27 & Shadow & 1044.17 & 1272.27 & 2281.08 & 367.79 \\
Bank28 & Shadow & 3402.82 & 2351.80 & 3232.02 & 447.19 \\
Bank29 & Shadow & 1132.35 & 862.18 & 2359.27 & 338.12 \\
\end{longtable}

\begin{longtable}{@{}lrrrrr@{}}
\caption{Shared initial data: interbank borrowing, wholesale funding, interbank assets, interbank liabilities, and project amount}\label{tab:paired-current-initial-balance-sheet-b}\\
\toprule
Bank & IB borrow. & Wholesale funding & IB assets & IB liab. & Project amount \\
\midrule
\endfirsthead
\multicolumn{6}{c}{\tablename\ \thetable\ -- continued}\\
\toprule
Bank & IB borrow. & Wholesale funding & IB assets & IB liab. & Project amount \\
\midrule
\endhead
\midrule
\multicolumn{6}{r}{Continued on the next page}\\
\endfoot
\bottomrule
\endlastfoot
CentralBank & 0.00 & 790.72 & 0.00 & 0.00 & 3943.08 \\
Bank1 & 0.00 & 754.74 & 0.00 & 0.00 & 3229.19 \\
Bank2 & 0.00 & 1118.05 & 0.00 & 0.00 & 3782.23 \\
Bank3 & 0.00 & 1304.78 & 0.00 & 0.00 & 4543.01 \\
Bank4 & 0.00 & 379.10 & 0.00 & 0.00 & 1109.74 \\
Bank5 & 0.00 & 757.99 & 0.00 & 0.00 & 2377.22 \\
Bank6 & 0.00 & 941.44 & 0.00 & 0.00 & 2950.86 \\
Bank7 & 0.00 & 860.99 & 0.00 & 0.00 & 2966.58 \\
Bank8 & 0.00 & 1057.51 & 0.00 & 0.00 & 3591.19 \\
Bank9 & 0.00 & 1142.47 & 0.00 & 0.00 & 4043.71 \\
Bank10 & 0.00 & 668.22 & 0.00 & 0.00 & 2187.29 \\
Bank11 & 0.00 & 534.35 & 0.00 & 0.00 & 1899.27 \\
Bank12 & 0.00 & 1212.90 & 0.00 & 0.00 & 3663.59 \\
Bank13 & 0.00 & 904.02 & 0.00 & 0.00 & 2818.56 \\
Bank14 & 0.00 & 1115.39 & 0.00 & 0.00 & 3777.91 \\
Bank15 & 0.00 & 1164.15 & 0.00 & 0.00 & 3267.06 \\
Bank16 & 0.00 & 760.44 & 0.00 & 0.00 & 3254.86 \\
Bank17 & 0.00 & 560.09 & 0.00 & 0.00 & 2062.06 \\
Bank18 & 0.00 & 325.36 & 0.00 & 0.00 & 1150.44 \\
Bank19 & 0.00 & 673.43 & 0.00 & 0.00 & 1916.53 \\
Bank20 & 0.00 & 642.63 & 0.00 & 0.00 & 1936.38 \\
Bank21 & 0.00 & 2575.39 & 0.00 & 0.00 & 2740.73 \\
Bank22 & 0.00 & 1708.35 & 0.00 & 0.00 & 2310.57 \\
Bank23 & 0.00 & 1211.93 & 0.00 & 0.00 & 1288.20 \\
Bank24 & 0.00 & 1996.17 & 0.00 & 0.00 & 2123.69 \\
Bank25 & 0.00 & 771.32 & 0.00 & 0.00 & 809.90 \\
Bank26 & 0.00 & 2667.03 & 0.00 & 0.00 & 2829.98 \\
Bank27 & 0.00 & 1913.29 & 0.00 & 0.00 & 2052.97 \\
Bank28 & 0.00 & 2784.83 & 0.00 & 0.00 & 4283.04 \\
Bank29 & 0.00 & 2021.15 & 0.00 & 0.00 & 2629.44 \\
\end{longtable}

\begin{longtable}{@{}lcrrrr@{}}
\caption{Shared initial data: activity, solvency, CAR, LCR, and leverage}\label{tab:paired-current-initial-state}\\
\toprule
Bank & Active & Solvency & CAR & LCR & Leverage \\
\midrule
\endfirsthead
\multicolumn{6}{c}{\tablename\ \thetable\ -- continued}\\
\toprule
Bank & Active & Solvency & CAR & LCR & Leverage \\
\midrule
\endhead
\midrule
\multicolumn{6}{r}{Continued on the next page}\\
\endfoot
\bottomrule
\endlastfoot
CentralBank & Yes & 1.4010 & 1.3810 & 3.0000 & 1.0105 \\
Bank1 & Yes & 0.5113 & 0.5080 & 1.2623 & 1.0126 \\
Bank2 & Yes & 0.1448 & 0.1609 & 0.6120 & 0.5915 \\
Bank3 & Yes & 0.6564 & 0.6274 & 1.5256 & 1.0757 \\
Bank4 & Yes & 0.6128 & 0.6809 & 1.7821 & 0.8597 \\
Bank5 & Yes & 0.6769 & 0.7521 & 1.9422 & 0.8713 \\
Bank6 & Yes & 0.5322 & 0.5913 & 1.5805 & 0.8418 \\
Bank7 & Yes & 0.5619 & 0.6243 & 1.6546 & 0.8489 \\
Bank8 & Yes & 0.1806 & 0.2007 & 0.7016 & 0.6437 \\
Bank9 & Yes & 0.6079 & 0.5986 & 1.4809 & 1.0263 \\
Bank10 & Yes & 0.8221 & 0.9135 & 2.3053 & 0.8916 \\
Bank11 & Yes & 0.9032 & 1.0035 & 2.5080 & 0.9003 \\
Bank12 & Yes & 0.5071 & 0.5634 & 1.5177 & 0.8353 \\
Bank13 & Yes & 0.2969 & 0.3299 & 0.9923 & 0.7481 \\
Bank14 & Yes & 0.3553 & 0.3948 & 1.1382 & 0.7804 \\
Bank15 & Yes & 0.3106 & 0.3237 & 0.8773 & 0.8852 \\
Bank16 & Yes & 0.2864 & 0.3182 & 0.9659 & 0.7412 \\
Bank17 & Yes & 0.6598 & 0.7331 & 1.8995 & 0.8684 \\
Bank18 & Yes & 0.4713 & 0.5237 & 1.4283 & 0.8250 \\
Bank19 & Yes & 0.4325 & 0.4805 & 1.3311 & 0.8122 \\
Bank20 & Yes & 0.3099 & 0.3444 & 1.0248 & 0.7561 \\
Bank21 & Yes & 0.5618 & 0.6242 & 1.6544 & 0.8489 \\
Bank22 & Yes & 0.5501 & 0.4833 & 1.0299 & 1.3353 \\
Bank23 & Yes & 0.4768 & 0.5298 & 1.4420 & 0.8266 \\
Bank24 & Yes & 0.3165 & 0.3517 & 1.0413 & 0.7599 \\
Bank25 & Yes & 2.1629 & 1.5000 & 3.0000 & 0.9558 \\
Bank26 & Yes & 0.3438 & 0.3820 & 1.1094 & 0.7747 \\
Bank27 & Yes & 0.4578 & 0.5086 & 1.3944 & 0.8207 \\
Bank28 & Yes & 1.0528 & 0.7945 & 1.8191 & 1.4469 \\
Bank29 & Yes & 0.4800 & 0.4306 & 0.9136 & 1.3134 \\
\end{longtable}

\begin{longtable}{@{}lrrrrr@{}}
\caption{Shared initial data: risk appetite, volatility, loan rate, investment rate, and outflow}\label{tab:paired-current-initial-state-b}\\
\toprule
Bank & Risk appetite & Volatility & Loan rate & Invest.\ rate & Outflow \\
\midrule
\endfirsthead
\multicolumn{6}{c}{\tablename\ \thetable\ -- continued}\\
\toprule
Bank & Risk appetite & Volatility & Loan rate & Invest.\ rate & Outflow \\
\midrule
\endhead
\midrule
\multicolumn{6}{r}{Continued on the next page}\\
\endfoot
\bottomrule
\endlastfoot
CentralBank & 0.3000 & 0.2391 & 0.000154 & 0.000216 & 0.2000 \\
Bank1 & 0.9825 & 0.3968 & 0.000153 & 0.000212 & 0.3542 \\
Bank2 & 0.9503 & 0.2153 & 0.000174 & 0.000214 & 0.3917 \\
Bank3 & 0.7598 & 0.3563 & 0.000175 & 0.000237 & 0.3763 \\
Bank4 & 0.9142 & 0.2115 & 0.000146 & 0.000207 & 0.4884 \\
Bank5 & 0.7265 & 0.3691 & 0.000152 & 0.000202 & 0.3107 \\
Bank6 & 0.7484 & 0.3102 & 0.000155 & 0.000227 & 0.3832 \\
Bank7 & 0.7738 & 0.3143 & 0.000178 & 0.000208 & 0.4677 \\
Bank8 & 0.8886 & 0.2965 & 0.000130 & 0.000220 & 0.4965 \\
Bank9 & 0.9556 & 0.3859 & 0.000165 & 0.000194 & 0.4617 \\
Bank10 & 0.8910 & 0.3698 & 0.000171 & 0.000230 & 0.4245 \\
Bank11 & 0.7375 & 0.3012 & 0.000131 & 0.000195 & 0.4838 \\
Bank12 & 0.8397 & 0.2365 & 0.000132 & 0.000221 & 0.4877 \\
Bank13 & 0.8399 & 0.2471 & 0.000172 & 0.000222 & 0.3134 \\
Bank14 & 0.7036 & 0.3090 & 0.000149 & 0.000216 & 0.4465 \\
Bank15 & 0.7404 & 0.3485 & 0.000154 & 0.000236 & 0.4459 \\
Bank16 & 0.7333 & 0.3820 & 0.000159 & 0.000223 & 0.3632 \\
Bank17 & 0.8151 & 0.3479 & 0.000179 & 0.000222 & 0.4992 \\
Bank18 & 0.9725 & 0.2305 & 0.000169 & 0.000206 & 0.3614 \\
Bank19 & 0.8984 & 0.2003 & 0.000168 & 0.000211 & 0.3552 \\
Bank20 & 0.8237 & 0.2880 & 0.000176 & 0.000231 & 0.3620 \\
Bank21 & 0.9566 & 0.3280 & 0.000178 & 0.000195 & 0.3423 \\
Bank22 & 0.9336 & 0.2572 & 0.000155 & 0.000234 & 0.4783 \\
Bank23 & 0.9884 & 0.3181 & 0.000153 & 0.000214 & 0.3859 \\
Bank24 & 0.9194 & 0.3761 & 0.000179 & 0.000220 & 0.4323 \\
Bank25 & 0.8190 & 0.3279 & 0.000143 & 0.000214 & 0.3942 \\
Bank26 & 0.9711 & 0.2327 & 0.000146 & 0.000233 & 0.4300 \\
Bank27 & 0.9540 & 0.2420 & 0.000148 & 0.000230 & 0.4759 \\
Bank28 & 0.8514 & 0.3591 & 0.000149 & 0.000210 & 0.3751 \\
Bank29 & 0.9661 & 0.3580 & 0.000141 & 0.000235 & 0.4312 \\
\end{longtable}

\endgroup
\endinput

\section{Archived GitHub Central-Policy File (Not Used in the Formal Sample)}

The following tables reproduce a separately supplied GitHub
central-policy initialization. The values differ from the shared current
file and were not used to generate the Chapter~7 results. They are
retained only for completeness of the data archive.

\begin{longtable}{@{}llrrrr@{}}
\caption{Archived GitHub file: core capital, liquid assets, current liabilities, and deposits}\label{tab:github-policy-initial-balance-sheet}\\
\toprule
Bank & Type & Core capital & Liquid assets & Current liab. & Deposits \\
\midrule
\endfirsthead
\multicolumn{6}{c}{\tablename\ \thetable\ -- continued}\\
\toprule
Bank & Type & Core capital & Liquid assets & Current liab. & Deposits \\
\midrule
\endhead
\midrule
\multicolumn{6}{r}{Continued on the next page}\\
\endfoot
\bottomrule
\endlastfoot
CentralBank & Central & 10000.00 & 5355.62 & 3556.24 & 1793.24 \\
Bank1 & Commercial & 4289.97 & 1367.79 & 3308.03 & 1899.83 \\
Bank2 & Commercial & 2095.70 & 1001.65 & 1583.03 & 863.87 \\
Bank3 & Commercial & 1968.23 & 1651.17 & 1378.33 & 875.07 \\
Bank4 & Commercial & 2304.89 & 1611.19 & 1882.25 & 1105.66 \\
Bank5 & Commercial & 2399.37 & 952.63 & 1944.78 & 1057.65 \\
Bank6 & Commercial & 1782.55 & 1433.28 & 1406.80 & 905.24 \\
Bank7 & Commercial & 4205.40 & 1542.07 & 3280.67 & 1952.65 \\
Bank8 & Commercial & 3773.50 & 1829.05 & 3372.01 & 2135.63 \\
Bank9 & Commercial & 1061.25 & 1256.85 & 758.17 & 443.58 \\
Bank10 & Commercial & 2929.86 & 1554.94 & 1889.89 & 1152.48 \\
Bank11 & Commercial & 2298.24 & 935.59 & 1867.02 & 1022.39 \\
Bank12 & Commercial & 2925.16 & 1392.40 & 2485.75 & 1618.73 \\
Bank13 & Commercial & 2064.39 & 1805.49 & 1725.42 & 1130.30 \\
Bank14 & Commercial & 3222.21 & 1717.33 & 2151.25 & 1396.70 \\
Bank15 & Commercial & 2065.61 & 1678.59 & 1536.21 & 846.45 \\
Bank16 & Commercial & 3105.53 & 1111.79 & 2245.30 & 1546.95 \\
Bank17 & Commercial & 2876.35 & 1172.35 & 2140.08 & 1177.52 \\
Bank18 & Commercial & 2228.04 & 708.43 & 1741.76 & 1090.79 \\
Bank19 & Commercial & 1692.92 & 1580.55 & 1331.99 & 751.99 \\
Bank20 & Commercial & 4030.56 & 1896.69 & 2921.97 & 1691.01 \\
Bank21 & Shadow & 988.26 & 1098.17 & 763.94 & 14.33 \\
Bank22 & Shadow & 4206.11 & 1437.06 & 2824.03 & 20.93 \\
Bank23 & Shadow & 2546.22 & 1720.64 & 1901.54 & 50.82 \\
Bank24 & Shadow & 4209.17 & 1097.50 & 2627.47 & 48.11 \\
Bank25 & Shadow & 2001.28 & 1802.95 & 1447.84 & 24.59 \\
Bank26 & Shadow & 3473.29 & 1333.14 & 2400.64 & 48.81 \\
Bank27 & Shadow & 3681.06 & 1750.88 & 2372.55 & 42.16 \\
Bank28 & Shadow & 1561.76 & 1736.00 & 1109.03 & 30.23 \\
Bank29 & Shadow & 1141.56 & 612.80 & 674.77 & 0.11 \\
\end{longtable}

\begin{longtable}{@{}lrrrrr@{}}
\caption{Archived GitHub file: interbank borrowing, wholesale funding, interbank assets, interbank liabilities, and project amount}\label{tab:github-policy-initial-balance-sheet-b}\\
\toprule
Bank & IB borrow. & Wholesale funding & IB assets & IB liab. & Project amount \\
\midrule
\endfirsthead
\multicolumn{6}{c}{\tablename\ \thetable\ -- continued}\\
\toprule
Bank & IB borrow. & Wholesale funding & IB assets & IB liab. & Project amount \\
\midrule
\endhead
\midrule
\multicolumn{6}{r}{Continued on the next page}\\
\endfoot
\bottomrule
\endlastfoot
CentralBank & 1126.34 & 636.66 & 0.00 & 0.00 & 3200.62 \\
Bank1 & 838.88 & 569.33 & 0.00 & 455.26 & 2977.23 \\
Bank2 & 326.45 & 392.70 & 0.00 & 0.00 & 1424.72 \\
Bank3 & 300.81 & 202.45 & 0.00 & 0.00 & 1240.50 \\
Bank4 & 383.95 & 392.65 & 0.00 & 0.00 & 1694.03 \\
Bank5 & 481.46 & 405.67 & 0.00 & 0.00 & 1750.30 \\
Bank6 & 222.75 & 278.82 & 0.00 & 0.00 & 1266.12 \\
Bank7 & 587.48 & 740.54 & 0.00 & 0.00 & 2952.60 \\
Bank8 & 792.74 & 443.64 & 0.00 & 0.00 & 3034.81 \\
Bank9 & 131.00 & 183.58 & 0.00 & 0.00 & 682.35 \\
Bank10 & 363.97 & 373.43 & 0.00 & 0.00 & 1700.90 \\
Bank11 & 467.92 & 376.71 & 190.41 & 0.00 & 1680.32 \\
Bank12 & 417.42 & 449.60 & 0.00 & 0.00 & 2237.17 \\
Bank13 & 224.58 & 370.53 & 0.00 & 0.00 & 1552.88 \\
Bank14 & 352.50 & 402.04 & 0.00 & 0.00 & 1936.12 \\
Bank15 & 385.80 & 303.96 & 0.00 & 0.00 & 1382.59 \\
Bank16 & 299.02 & 399.34 & 444.27 & 0.00 & 2020.77 \\
Bank17 & 556.58 & 405.99 & 0.00 & 0.00 & 1926.08 \\
Bank18 & 363.29 & 287.68 & 0.00 & 52.18 & 1567.59 \\
Bank19 & 336.73 & 243.27 & 0.00 & 0.00 & 1198.79 \\
Bank20 & 670.10 & 560.86 & 0.00 & 0.00 & 2629.77 \\
Bank21 & 284.92 & 464.69 & 0.00 & 0.00 & 687.55 \\
Bank22 & 1035.47 & 1767.62 & 0.00 & 0.00 & 2541.63 \\
Bank23 & 650.38 & 1200.34 & 0.00 & 0.00 & 1711.39 \\
Bank24 & 886.34 & 1693.02 & 0.00 & 317.19 & 2364.73 \\
Bank25 & 527.95 & 895.30 & 0.00 & 0.00 & 1303.06 \\
Bank26 & 805.68 & 1546.15 & 189.95 & 0.00 & 2160.58 \\
Bank27 & 821.96 & 1508.43 & 0.00 & 0.00 & 2135.30 \\
Bank28 & 365.75 & 713.05 & 0.00 & 0.00 & 998.13 \\
Bank29 & 217.74 & 456.92 & 0.00 & 0.00 & 607.29 \\
\end{longtable}

\begin{longtable}{@{}lcrrrr@{}}
\caption{Archived GitHub file: activity, solvency, CAR, LCR, and leverage}\label{tab:github-policy-initial-state}\\
\toprule
Bank & Active & Solvency & CAR & LCR & Leverage \\
\midrule
\endfirsthead
\multicolumn{6}{c}{\tablename\ \thetable\ -- continued}\\
\toprule
Bank & Active & Solvency & CAR & LCR & Leverage \\
\midrule
\endhead
\midrule
\multicolumn{6}{r}{Continued on the next page}\\
\endfoot
\bottomrule
\endlastfoot
CentralBank & Yes & 4.3179 & 3.1244 & 7.5299 & 1.8672 \\
Bank1 & Yes & 1.5727 & 1.4409 & 0.6896 & 3.1364 \\
Bank2 & Yes & 1.9566 & 1.4709 & 1.5819 & 2.0923 \\
Bank3 & Yes & 2.6259 & 1.5866 & 2.9949 & 1.1920 \\
Bank4 & Yes & 2.0805 & 1.3606 & 2.1400 & 1.4306 \\
Bank5 & Yes & 1.7236 & 1.3708 & 1.2246 & 2.5187 \\
Bank6 & Yes & 2.2859 & 1.4079 & 2.5471 & 1.2437 \\
Bank7 & Yes & 1.7519 & 1.4243 & 1.1751 & 2.7271 \\
Bank8 & Yes & 1.6615 & 1.2434 & 1.3561 & 2.0631 \\
Bank9 & Yes & 3.0575 & 1.5553 & 4.1444 & 0.8444 \\
Bank10 & Yes & 2.3731 & 1.7225 & 2.0569 & 1.8842 \\
Bank11 & Yes & 1.8341 & 1.2944 & 1.5077 & 2.0411 \\
Bank12 & Yes & 1.7369 & 1.3075 & 1.4004 & 2.1008 \\
Bank13 & Yes & 2.2429 & 1.3294 & 2.6160 & 1.1434 \\
Bank14 & Yes & 2.2961 & 1.6643 & 1.9957 & 1.8763 \\
Bank15 & Yes & 2.4373 & 1.4940 & 2.7317 & 1.2306 \\
Bank16 & Yes & 2.0762 & 1.3846 & 1.7326 & 1.9958 \\
Bank17 & Yes & 1.8918 & 1.4934 & 1.3695 & 2.4535 \\
Bank18 & Yes & 1.6560 & 1.4213 & 0.9419 & 3.1450 \\
Bank19 & Yes & 2.4576 & 1.4122 & 2.9665 & 1.0711 \\
Bank20 & Yes & 2.0285 & 1.5327 & 1.6228 & 2.1251 \\
Bank21 & Yes & 2.7311 & 1.4374 & 3.5938 & 0.8999 \\
Bank22 & Yes & 1.9983 & 1.6549 & 1.2722 & 2.9269 \\
Bank23 & Yes & 2.2439 & 1.4878 & 2.2622 & 1.4798 \\
Bank24 & Yes & 1.8990 & 1.7800 & 0.7425 & 3.8352 \\
Bank25 & Yes & 2.6275 & 1.5358 & 3.1132 & 1.1100 \\
Bank26 & Yes & 2.0813 & 1.5399 & 1.5861 & 2.2804 \\
Bank27 & Yes & 2.2895 & 1.7239 & 1.8449 & 2.1024 \\
Bank28 & Yes & 2.9736 & 1.5647 & 3.9133 & 0.8996 \\
Bank29 & Yes & 2.6000 & 1.8798 & 2.2704 & 1.8629 \\
\end{longtable}

\begin{longtable}{@{}lrrrrr@{}}
\caption{Archived GitHub file: risk appetite, volatility, loan rate, investment rate, and outflow}\label{tab:github-policy-initial-state-b}\\
\toprule
Bank & Risk appetite & Volatility & Loan rate & Invest.\ rate & Outflow \\
\midrule
\endfirsthead
\multicolumn{6}{c}{\tablename\ \thetable\ -- continued}\\
\toprule
Bank & Risk appetite & Volatility & Loan rate & Invest.\ rate & Outflow \\
\midrule
\endhead
\midrule
\multicolumn{6}{r}{Continued on the next page}\\
\endfoot
\bottomrule
\endlastfoot
CentralBank & 0.3000 & 0.8246 & 0.000314 & 0.000307 & 0.2000 \\
Bank1 & 0.2726 & 0.7815 & 0.000267 & 0.000283 & 0.3917 \\
Bank2 & 0.1728 & 0.8041 & 0.000298 & 0.000299 & 0.4784 \\
Bank3 & 0.0686 & 0.7080 & 0.000275 & 0.000293 & 0.4639 \\
Bank4 & 0.2622 & 0.6701 & 0.000311 & 0.000303 & 0.4733 \\
Bank5 & 0.2951 & 0.8054 & 0.000308 & 0.000286 & 0.4813 \\
Bank6 & 0.1172 & 0.6208 & 0.000306 & 0.000295 & 0.4183 \\
Bank7 & 0.0166 & 0.9944 & 0.000277 & 0.000295 & 0.3899 \\
Bank8 & 0.2153 & 0.8579 & 0.000273 & 0.000297 & 0.4610 \\
Bank9 & 0.0335 & 0.8511 & 0.000313 & 0.000324 & 0.3679 \\
Bank10 & 0.1140 & 0.6325 & 0.000304 & 0.000329 & 0.4056 \\
Bank11 & 0.2746 & 0.6911 & 0.000262 & 0.000302 & 0.4594 \\
Bank12 & 0.0767 & 0.7833 & 0.000274 & 0.000302 & 0.4652 \\
Bank13 & 0.0408 & 0.7533 & 0.000297 & 0.000302 & 0.4224 \\
Bank14 & 0.1169 & 0.7726 & 0.000267 & 0.000293 & 0.3338 \\
Bank15 & 0.2265 & 0.9998 & 0.000275 & 0.000320 & 0.3048 \\
Bank16 & 0.0099 & 0.7210 & 0.000261 & 0.000316 & 0.3654 \\
Bank17 & 0.2767 & 0.9675 & 0.000264 & 0.000298 & 0.4138 \\
Bank18 & 0.0765 & 0.8197 & 0.000263 & 0.000290 & 0.4006 \\
Bank19 & 0.2561 & 0.8479 & 0.000280 & 0.000303 & 0.3257 \\
Bank20 & 0.1179 & 0.8603 & 0.000308 & 0.000310 & 0.4222 \\
Bank21 & 0.1527 & 0.6293 & 0.000307 & 0.000298 & 0.3359 \\
Bank22 & 0.0524 & 0.7790 & 0.000271 & 0.000296 & 0.4254 \\
Bank23 & 0.2101 & 0.7341 & 0.000290 & 0.000283 & 0.3301 \\
Bank24 & 0.1203 & 0.8892 & 0.000304 & 0.000300 & 0.4081 \\
Bank25 & 0.1293 & 0.7952 & 0.000308 & 0.000286 & 0.3708 \\
Bank26 & 0.1479 & 0.9776 & 0.000264 & 0.000303 & 0.3559 \\
Bank27 & 0.1206 & 0.7941 & 0.000268 & 0.000299 & 0.3467 \\
Bank28 & 0.2038 & 0.6607 & 0.000281 & 0.000326 & 0.4902 \\
Bank29 & 0.0010 & 0.9008 & 0.000312 & 0.000297 & 0.4269 \\
\end{longtable}

